% Allgemeine endliche Elternsysteme, nach Wurzelwegen sowie Eltern und Kindern.
\section{Endliche Wurzelbäume und gerichtete Elternkanten}

\providecommand{\BXXVIIFiniteRooted}[3]{%
  \mathsf{EndWBaum}(#1,#2,#3)}
\providecommand{\BXXVIIParentSystem}[4]{%
  \mathsf{ElterSys}(#1,#2,#3,#4)}
\providecommand{\BXXVIISymmetricEdges}[2]{%
  \operatorname{sym}_{#1}(#2)}

\noindent\begin{minipage}{\linewidth}
\FormulaDefDeltaK[Endlicher verwurzelter Baum]{%
  \BXXVIIFiniteRooted VEr\coloneqq
  \Finite V\land\RootedTreeStruct{V,E,r}
}{B27FiniteRootedTreeDef}{%
  \DeltaRow{Mengen}{V\dsep E\dsep r}
  \DeltaRow{Neue Struktur}{\BXXVIIFiniteRooted VEr}
}
\end{minipage}\par

Die Endlichkeit ist genau die Endlichkeit des Knotenträgers aus Band~20.
Eltern und Kinder behalten die zuvor eingeführten Definitionen. Eine Beschriftung ist eine zusätzliche Abbildung auf den
betreffenden Knoten; sie verändert weder deren Identität noch die Kanten.

\noindent\begin{minipage}{\linewidth}
\FormulaThmDeltaK[Die Kindermengen eines endlichen Baums sind endlich]{%
  \BXXVIIFiniteRooted VEr\dsep a\in V\vdash
  \Finite{\TreeChildren VEr a}
}{B27FiniteRootedChildrenFinite}{%
  \DeltaRow{Mengen}{V\dsep E\dsep r\dsep a\dsep x}
}
\end{minipage}\par
\begin{tabproofwide}
  \proofstepwidestar[1]{\BXXVIIFiniteRooted VEr}{\rA}
  \proofstepwidestar[2]{a\in V}{\rA}
  \proofstepwidestar[1]{\Finite V\land\RootedTreeStruct{V,E,r}}{%
    \FormulaRefAuto{B27FiniteRootedTreeDef}[def]{1}}
  \proofstepwidestar[4]{x\in\TreeChildren VEr a}{\rA}
  \proofstepwidestar[1,2]{x\in\TreeChildren VEr a\leftrightarrow
    \TreeParentRel VEr a x}{%
    \FormulaRefAuto{TreeChildrenMembership}[thm]{\rAEb{3},2}}
  \proofstepwidestar[1,2,4]{\TreeParentRel VEr a x}{%
    \rRE{4,\rLREa{5}}}
  \proofstepwidestar[1,2,4]{x\in V\setminus\{r\}}{%
    \FormulaRefAuto{TreeParentAccess}[thm]{\rAEb{3},6}}
  \proofstepwidestar[1,2,4]{x\in V}{%
    \FormulaRefAuto{x\in A\setminus B\vdash x\in A}[thm]{7}}
  \proofstepwidestar[1,2]{\TreeChildren VEr a\subseteq V}{%
    \FormulaRefAuto{SubsetIntroductionRule}[thm]{[4]\vdots8}}
  \proofstepwidestar[1,2]{\Finite{\TreeChildren VEr a}}{%
    \FormulaRefAuto{FiniteSubset}[thm]{9,\rAEa{3}}}
\end{tabproofwide}

\subsection{Ein gerichtetes Einführungskriterium}

Eine gerichtete Kante \((a,x)\) soll von einem Elternknoten zu seinem Kind
zeigen. Eine natürliche Höhenmarke verhindert, dass man durch solche
Kanten zu einem früheren Knoten zurückkehrt. Die Marken müssen nicht die
Abstände zur Wurzel sein; es genügt, dass sie an jeder Kante streng wachsen.

\noindent\begin{minipage}{\linewidth}
\FormulaDefDeltaK[Endliches Elternsystem mit Höhenmarken]{%
  \begin{aligned}[t]
    \BXXVIIParentSystem VDrh\coloneqq{}&
      \Finite V\land D\subseteq V\times V\land r\in V\\[-2pt]
    &\land h\colon V\to\mathbb N\land h(r)=0\\[-2pt]
    &\land\bigl(\forall x\in V\setminus\{r\}\,
       \exists!a\in V\,((a,x)\in D)\bigr)\\[-2pt]
    &\land\forall a,x\in V\,
       ((a,x)\in D\rightarrow h(a)<h(x)).
  \end{aligned}
}{B27RankedParentSystemDef}{%
  \DeltaRow{Mengen}{V\dsep D\dsep r\dsep h\dsep a\dsep x}
  \DeltaRow{Neue Struktur}{\BXXVIIParentSystem VDrh}
}
\end{minipage}\par

\noindent\begin{minipage}{\linewidth}
\FormulaDefDeltaK[Vergessen der Kantenrichtung]{%
  \BXXVIISymmetricEdges VD\coloneqq
  \{(a,x)\in V\times V\mid(a,x)\in D\lor(x,a)\in D\}
}{B27SymmetricParentEdgesDef}{%
  \DeltaRow{Mengen}{V\dsep D\dsep a\dsep x}
  \DeltaRow{Neue Menge}{\BXXVIISymmetricEdges VD}
}
\end{minipage}\par

% Zitierbare Zwischenergebnisse des allgemeinen Elternsystem-Kriteriums.
\FormulaThmDeltaK[Graphdaten eines Elternsystems]{%
  \begin{aligned}[t]
    \BXXVIIParentSystem VDrh\vdash{}&\GraphStruct{V,D}\land
      \SimpleUndirectedGraphStruct{V,\BXXVIISymmetricEdges VD}\land{}\\[-2pt]
    &r\in V\land\forall a\in V\;((a,r)\notin D)
  \end{aligned}
}{B27RankedParentGraphData}{\DeltaRow{Mengen}{V\dsep D\dsep r\dsep h\dsep a\dsep x}}
\begin{proof}
\emergencystretch=2em
\FormulaRefAuto{B27RankedParentSystemDef}[def] liefert
\(D\subseteq V\times V\), \(r\in V\), \(h\colon V\to\mathbb N\),
\(h(r)=0\) und die strikte Höhenzunahme an jeder Kante.
Nach \FormulaRefAuto{GraphDef}[def] ist \((V,D)\) ein Graph.
Aus \FormulaRefAuto{B27SymmetricParentEdgesDef}[def] und
\FormulaRefAuto{ProductComprehensionMembership}[thm] folgt für
\(a,x\in V\)
\[
 (a,x)\in\BXXVIISymmetricEdges VD
 \leftrightarrow ((a,x)\in D\lor(x,a)\in D).
\]
Die definierte Relation liegt in \(V\times V\); die rechte Seite ist
unter Vertauschung von \(a,x\) invariant. Daher gilt
\FormulaRefAuto{UndirectedGraphIntroduction}[thm]. Eine Schleife würde
in beiden Disjunktionsfällen \(h(x)<h(x)\) ergeben. Dies widerspricht
\FormulaRefAuto{PeanoStrictEqualityImpossible}[thm] mit der reflexiven
Gleichheit \(h(x)=h(x)\). Mit
\FormulaRefAuto{SimpleUndirectedGraphIntroduction}[thm] folgt die
Schlichtheit. Für \(a\in V\) würde \((a,r)\in D\) entsprechend
\(h(a)<h(r)=0\) liefern. Der Funktionswert liegt in \(\mathbb N\)
nach \FormulaRefAuto{F\colon A\to B\dsep x\in A\vdash F(x)\in B}[thm];
\FormulaRefAuto{PeanoNoLtZero}[thm] schließt die Kante aus.
\end{proof}

\FormulaThmDeltaK[Höhenzunahme entlang eines gerichteten Elternwegs]{%
  \begin{aligned}[t]
    &\BXXVIIParentSystem VDrh\dsep\DirectedPath VDp n\dsep
      i,j\in\NatSeg{n+1}\dsep i<j\\[-2pt]
    &\vdash h(p(i))<h(p(j))
  \end{aligned}
}{B27RankedPathHeight}{\DeltaRow{Mengen}{V\dsep D\dsep r\dsep h\dsep p\dsep n\dsep i\dsep j}}
\begin{proof}
\emergencystretch=2em
Die Typdaten von \(p\) folgen aus
\FormulaRefAuto{B27RankedParentGraphData}[thm] und
\FormulaRefAuto{FinitePathTyping}[thm]. Für jeden Schrittindex
\(k\in\NatSeg n\) liefern
\FormulaRefAuto{DirectedPathWalkAxiom}[ax] und
\FormulaRefAuto{DirectedWalkStepEdge}[thm] die Kante
\((p(k),p(k+1))\in D\). Nach
\FormulaRefAuto{B27RankedParentSystemDef}[def] gilt deshalb
\(h(p(k))<h(p(k+1))\).

Für die festen Parameter \(i,n\) verwenden wir das Prädikat
\[
 Q(d)\coloneqq\bigl(i+(d+1)\in\NatSeg{n+1}\rightarrow
   h(p(i))<h(p(i+(d+1)))\bigr).
\]
Induktion über \(d\) mit
\FormulaRefAuto{PeanoInductionAddOne}[thm] verkettet die Schrittungleichungen.
Bei \(d=0\) ist die Behauptung unter ihrer Indexvoraussetzung die
soeben bewiesene Kantenungleichung. Für \(Q(d+1)\) nehmen wir
\(i+((d+1)+1)\in\NatSeg{n+1}\) an. Dann ist auch der Zwischenindex
\(i+(d+1)\) zulässig; die Anfangsabschnittskriterien
\FormulaRefAuto{PeanoNatSegStrictMembership}[thm] und
\FormulaRefAuto{FinitePathPrefixIndexData}[thm] sichern alle benötigten
Indexzugehörigkeiten. Die Induktionsungleichung und die letzte
Schrittungleichung ergeben \(Q(d+1)\) durch
\FormulaRefAuto{PeanoLtTransitive}[thm]; die Schreibweisen der Indizes
stimmen nach \FormulaRefAuto{PeanoAddAssociative}[thm] überein.
Aus dem ursprünglich gegebenen \(i<j\) liefert
\FormulaRefAuto{PeanoStrictAddOneDistance}[thm] schließlich ein
\(d\in\mathbb N\) mit \(j=i+(d+1)\). Die nun bewiesene Aussage
\(Q(d)\) liefert wegen \(j\in\NatSeg{n+1}\) die Behauptung.
\end{proof}

\FormulaThmDeltaK[Gerichtete Erreichbarkeit vom Wurzelknoten]{%
  \BXXVIIParentSystem VDrh\dsep x\in V\vdash
    \exists z\;(z\in\GraphPathsBetween VDrx)
}{B27RankedDirectedReachability}{\DeltaRow{Mengen}{V\dsep D\dsep r\dsep h\dsep x\dsep n}}
\begin{proof}
\emergencystretch=2em
Wir beweisen für jedes \(n\in\mathbb N\): Jeder Knoten \(x\in V\)
mit \(h(x)\leq n\) ist durch einen gerichteten einfachen Weg erreichbar.
Das Induktionsprinzip ist \FormulaRefAuto{PeanoInductionAddOne}[thm].
Für \(n=0\) gilt \(h(x)=0\), denn
\FormulaRefAuto{PeanoZeroLeq}[thm] und
\FormulaRefAuto{PeanoLeqAntisymmetric}[thm] liefern beide Richtungen.
Wäre \(x\ne r\), so gäbe es nach
\FormulaRefAuto{B27RankedParentSystemDef}[def] einen Elternknoten
\(a\in V\) mit \(h(a)<h(x)=0\), entgegen
\FormulaRefAuto{PeanoNoLtZero}[thm]. Also ist \(x=r\).
Der Graph \(p=\{(0,r)\}\) ist nach
\FormulaRefAuto{SingletonGraphFunction}[thm],
\FormulaRefAuto{SingletonGraphValue}[thm] und
\FormulaRefAuto{PeanoNatSegOneEquality}[thm] eine injektive Folge auf
\(\NatSeg1\) mit Anfang und Ende \(r\). Die Kantenbedingung ist auf
\(\NatSeg0=\varnothing\) leer. Somit gelten
\FormulaRefAuto{DirectedPathCharacterization}[thm] und
\FormulaRefAuto{FiniteSimplePathFamilyMembership}[thm].

Im Induktionsschritt sei \(h(x)\leq n+1\). Falls \(h(x)\leq n\),
gilt die Induktionsannahme. Andernfalls liefert
\FormulaRefAuto{PeanoLeqSplit}[thm], zusammen mit
\FormulaRefAuto{PeanoLeqAddOneStrictBackward}[thm], die Gleichheit \(h(x)=n+1\).
Wegen \(h(r)=0\) und \FormulaRefAuto{PeanoZeroNeAddOne}[thm] ist
\(x\ne r\). Wähle den durch
\FormulaRefAuto{B27RankedParentSystemDef}[def] gegebenen Elternknoten
\(a\in V\). Aus \(h(a)<n+1\) folgt \(h(a)\leq n\) nach
\FormulaRefAuto{PeanoLeqAddOneStrictBackward}[thm]. Die Induktionsannahme und
\FormulaRefAuto{FiniteSimplePathFamilyMembership}[thm] liefern einen
einfachen Weg \(p\) mit \(p(0)=r\) und \(p(m)=a\).
Kein Wert von \(p\) kann \(x\) sein: Am Endindex wäre sonst
\(h(a)=h(x)\); an einem früheren Index ergäbe
\FormulaRefAuto{B27RankedPathHeight}[thm] sogar \(h(x)<h(a)\).
Im ersten Fall widerspricht dies \(h(a)<h(x)\) nach
\FormulaRefAuto{PeanoStrictEqualityImpossible}[thm]. Im zweiten Fall
liefert \FormulaRefAuto{PeanoLtTransitive}[thm] die unmögliche
Ungleichung \(h(x)<h(x)\).
Das Bildmengenkriterium
\FormulaRefAuto{F\colon A\to B\dsep y\in F[A]\vdash\exists x\in A\;y=F(x)}[thm]
liefert folglich \(x\notin p[\NatSeg{m+1}]\).
\FormulaRefAuto{FinitePathFreshEndpointExtension}[thm] hängt die Kante
\((a,x)\) an; \FormulaRefAuto{FiniteSimplePathFamilyMembership}[thm]
liefert das neue gespeicherte Wegobjekt. Damit ist der Induktionsschritt
bewiesen. Für den gegebenen Knoten setzen wir schließlich \(n=h(x)\);
Natürlichkeit und \FormulaRefAuto{PeanoLeqReflexive}[thm] sind erfüllt.
\end{proof}

\FormulaThmDeltaK[Jeder einfache Wurzelweg folgt der Elternrichtung]{%
  \begin{aligned}[t]
    &\BXXVIIParentSystem VDrh\dsep
      \DirectedPath V{\BXXVIISymmetricEdges VD}p n\dsep p(0)=r\vdash
      \DirectedPath VDp n
  \end{aligned}
}{B27RankedRootPathForward}{\DeltaRow{Mengen}{V\dsep D\dsep r\dsep h\dsep p\dsep n\dsep i}}
\begin{proof}
\emergencystretch=2em
Nach \FormulaRefAuto{FinitePathTyping}[thm] ist \(p\) eine injektive
Folge auf \(\NatSeg{n+1}\). Für das feste \(n\) setzen wir
\[
 Q(i)\coloneqq\bigl(i\in\NatSeg n\rightarrow(p(i),p(i+1))\in D\bigr).
\]
Wir beweisen \(Q(i)\) über den natürlichen Schrittindex; das Prinzip ist
\FormulaRefAuto{PeanoInductionAddOne}[thm]. Das Kantenkriterium aus
\FormulaRefAuto{B27SymmetricParentEdgesDef}[def] und
\FormulaRefAuto{ProductComprehensionMembership}[thm] liefert jeweils
die Vorwärts- oder Rückwärtskante.

Ist der jeweils betrachtete Index nicht in \(\NatSeg n\), ist
\(Q(i)\) wegen der falschen Voraussetzung erfüllt. Dies erfasst
insbesondere den Fall \(n=0\). Beim ersten zulässigen Schritt würde
die Rückwärtskante nach \(p(0)=r\)
führen. Dies widerspricht
\FormulaRefAuto{B27RankedParentGraphData}[thm]. Angenommen, der Schritt
von \(p(i)\) nach \(p(i+1)\) liegt in \(D\), und der nächste Schritt
läge rückwärts: \((p(i+2),p(i+1))\in D\).
Der Knoten \(p(i+1)\) ist von \(r=p(0)\) verschieden, da \(p\)
injektiv und \(i+1\ne0\) ist
(\FormulaRefAuto{PeanoZeroNeAddOne}[thm]). Die eindeutige
Elternexistenz in \FormulaRefAuto{B27RankedParentSystemDef}[def]
und \FormulaRefAuto{\exists! x\,P(x),\; P(a),\; P(b) \vdash a = b}[thm]
ergeben \(p(i)=p(i+2)\). Injektivität liefert \(i=i+2\).
Andererseits liefern zwei Anwendungen von
\FormulaRefAuto{PeanoLtAddOne}[thm] die Ungleichungen
\(i<i+1\) und \(i+1<i+2\). Nach
\FormulaRefAuto{PeanoLtTransitive}[thm] gilt also \(i<i+2\), im
Widerspruch zu \FormulaRefAuto{PeanoStrictEqualityImpossible}[thm].
Also ist auch der nächste
Schritt vorwärts gerichtet. Die Indexzugehörigkeiten in diesem Schritt
folgen aus \FormulaRefAuto{FinitePathPrefixIndexData}[thm] und
\FormulaRefAuto{FinitePathIndexShiftOne}[thm].
Die Folge und ihre Injektivität bleiben unverändert; die bewiesenen
Kantenbedingungen liefern mit
\FormulaRefAuto{DirectedPathCharacterization}[thm] den gerichteten Pfad.
\end{proof}

\Needspace{14\baselineskip}
\FormulaThmDeltaK[Eindeutigkeit gerichteter Wurzelwege]{%
  \begin{aligned}[t]
    &\BXXVIIParentSystem VDrh\dsep x\in V\dsep
      z,z'\in\GraphPathsBetween VDrx\vdash z=z'
  \end{aligned}
}{B27RankedDirectedPathUnique}{\DeltaRow{Mengen}{V\dsep D\dsep r\dsep h\dsep x\dsep z\dsep z'}}
\begin{proof}
\emergencystretch=2em
\FormulaRefAuto{FiniteSimplePathFamilyMembership}[thm] schreibt
\(z=(m,p)\), \(z'=(n,q)\) mit natürlichen Längen und einfachen
gerichteten Wegen, beide von \(r\) nach \(x\).
Für die festen Parameter \(V,D,r,h\) betrachten wir
\[
\begin{aligned}
 \mathcal C=\{s\in\mathbb N\mid{}&\exists x\in V\,\exists m,n\in\mathbb N\,
   \exists p,q\in\powerset(\mathbb N\times V)\\[-2pt]
 &((m,p),(n,q)\in\GraphPathsBetween VDrx\land{}\\[-2pt]
 &\quad s=m+n\land(m,p)\ne(n,q))\}.
\end{aligned}
\]
Der Endpunkt wird hier ausdrücklich mitquantifiziert. Die Menge ist
durch Aussonderung definiert. Wäre sie nicht leer, gäbe
\FormulaRefAuto{NaturalWellOrderingPrinciple}[thm] eine kleinste
Gegenbeispiellänge. Wähle zu ihr die angezeigten Endpunkt- und
Wegzeugen. Sie hätten folgende Eigenschaften.

Ist eine Länge null, so ist \(x=r\). Eine positive andere Länge
hätte eine letzte eingehende Kante nach \(r\), was
\FormulaRefAuto{B27RankedParentGraphData}[thm] ausschließt.
Also sind beide Längen null. Die Folgen sind dann auf
\(\NatSeg1=\{0\}\) gleich; dies folgt aus
\FormulaRefAuto{PeanoNatSegOneEquality}[thm] und
\FormulaRefAuto{FunctionExtensionality}[thm].

Sind beide Längen positiv, liefert
\FormulaRefAuto{PeanoNonzeroUniquePredecessor}[thm]
\(m=k+1\) und \(n=l+1\). Der gemeinsame Endpunkt ist wegen
\FormulaRefAuto{B27RankedParentGraphData}[thm] keine Wurzel.
Die beiden letzten Kanten geben daher dieselben Eltern:
\(p(k)=q(l)\), nach \FormulaRefAuto{B27RankedParentSystemDef}[def]
und der Eindeutigkeitsregel
\FormulaRefAuto{\exists! x\,P(x),\; P(a),\; P(b) \vdash a = b}[thm].
\FormulaRefAuto{FinitePathInitialSegment}[thm] liefert die beiden
verkürzten Wege \(p\restriction_{\NatSeg{k+1}}\) und
\(q\restriction_{\NatSeg{l+1}}\) zu diesem gemeinsamen Elternknoten.
Ihre Längensumme ist kleiner: \(k+l<m+n\) folgt aus
\FormulaRefAuto{PeanoLtAddOne}[thm] und
\FormulaRefAuto{PeanoAddStrictMonotoneBoth}[thm]. Die Minimalität
der gewählten Gegenbeispiellänge erzwingt deshalb die Gleichheit
der verkürzten Wegobjekte. Nach
\FormulaRefAuto{(a,b)=(c,d)\eqvdash a=c\land b=d}[thm]
sind \(k=l\) und die beiden Einschränkungen gleich.
Somit sind \(m=n\), alle früheren Koordinaten gleich und beide
letzten Koordinaten gleich \(x\). Die Aufteilung in den früheren
Anfangsabschnitt und den letzten Index folgt aus
\FormulaRefAuto{PeanoNatSegAddOneMembershipIff}[thm].
\FormulaRefAuto{FunctionExtensionality}[thm] ergibt \(p=q\),
und das Paargleichheitskriterium liefert \(z=z'\).
Es gibt daher kein kleinstes und folglich kein Gegenbeispiel.
\end{proof}

\FormulaThmDeltaK[Eindeutige Wege im symmetrisierten Elternsystem]{%
  \BXXVIIParentSystem VDrh\vdash
    \forall x\in V\,\exists!z\;(z\in\GraphPathsBetween V{\BXXVIISymmetricEdges VD}rx)
}{B27RankedSymmetricRootPaths}{\DeltaRow{Mengen}{V\dsep D\dsep r\dsep h\dsep x\dsep z}}
\begin{proof}
\emergencystretch=2em
Für \(x\in V\) liefert
\FormulaRefAuto{B27RankedDirectedReachability}[thm] einen gerichteten
Wurzelweg. Jede Kante aus \(D\) liegt nach
\FormulaRefAuto{B27SymmetricParentEdgesDef}[def] auch in der
symmetrisierten Relation. Folglich bleibt die gleiche injektive Folge
ein Pfad, nach \FormulaRefAuto{DirectedPathCharacterization}[thm].
Mit \FormulaRefAuto{FiniteSimplePathFamilyMembership}[thm] ist also
ein gespeicherter symmetrisierter Wurzelweg vorhanden.

Sind zwei solche Wegobjekte gegeben, so liefern
\FormulaRefAuto{FiniteSimplePathFamilyMembership}[thm] und
\FormulaRefAuto{B27RankedRootPathForward}[thm] gerichtete Wege
mit denselben Folgen und Längen. Ihre Gleichheit folgt aus
\FormulaRefAuto{B27RankedDirectedPathUnique}[thm]. Die Regel
\UEI{} verbindet den vorhandenen Zeugen mit dieser Eindeutigkeit;
\rUI{} schließt über den beliebigen Knoten \(x\).
\end{proof}


\noindent\begin{minipage}{\linewidth}
\FormulaThmDeltaK[Ein Elternsystem trägt einen endlichen Wurzelbaum]{%
  \BXXVIIParentSystem VDrh\vdash
    \BXXVIIFiniteRooted V{\BXXVIISymmetricEdges VD}r
}{B27RankedParentSystemTree}{%
  \DeltaRow{Mengen}{V\dsep D\dsep r\dsep h}
}
\end{minipage}\par

\begin{tabproofwide}
  \proofstepwidestar[1]{\BXXVIIParentSystem VDrh}{\rA}
  \proofstepwidestar[1]{\Finite V}{\FormulaRefAuto{B27RankedParentSystemDef}[def]{1}}
  \proofstepwidestar[1]{\begin{aligned}[t]
    &\GraphStruct{V,D}\land
      \SimpleUndirectedGraphStruct{V,\BXXVIISymmetricEdges VD}\land{}\\[-2pt]
    &r\in V\land\forall a\in V\;((a,r)\notin D)
  \end{aligned}}{\FormulaRefAuto{B27RankedParentGraphData}[thm]{1}}
  \proofstepwidestar[1]{\SimpleUndirectedGraphStruct{V,\BXXVIISymmetricEdges VD}}{\rAEa{\rAEb{3}}}
  \proofstepwidestar[1]{r\in V}{\rAEa{\rAEb{\rAEb{3}}}}
  \proofstepwidestar[1]{\forall x\in V\,\exists!z\;
    (z\in\GraphPathsBetween V{\BXXVIISymmetricEdges VD}rx)}{\FormulaRefAuto{B27RankedSymmetricRootPaths}[thm]{1}}
  \proofstepwidestar[1]{\RootedTreeStruct{V,\BXXVIISymmetricEdges VD,r}}{\FormulaRefAuto{RootedTreeFromUniqueRootPaths}[thm]{4,5,6}}
  \proofstepwidestar[1]{\BXXVIIFiniteRooted V{\BXXVIISymmetricEdges VD}r}{\FormulaRefAuto{B27FiniteRootedTreeDef}[def]{\rAI{2,7}}}
\end{tabproofwide}

\noindent\begin{minipage}{\linewidth}
\FormulaThmDeltaK[Die gerichteten Kanten sind genau die Baumeltern]{%
  \begin{aligned}[t]
    &\BXXVIIParentSystem VDrh\dsep a,x\in V\vdash\\[-2pt]
    &(a,x)\in D\leftrightarrow
      \TreeParentRel V{\BXXVIISymmetricEdges VD}r a x
  \end{aligned}
}{B27RankedParentEdgesAgree}{%
  \DeltaRow{Mengen}{V\dsep D\dsep r\dsep h\dsep a\dsep x}
}
\end{minipage}\par
\noindent\textit{Beweis.}
Nach \FormulaRefAuto{B27RankedParentSystemTree}[thm] ist die
zuvor definierte Elternrelation anwendbar. Gilt \((a,x)\in D\), so verlängert
diese Kante den nach \FormulaRefAuto{B27RankedDirectedReachability}[thm]
vorhandenen gerichteten Wurzelweg nach \(a\) zu einem einfachen
Wurzelweg nach \(x\): Nach \FormulaRefAuto{B27RankedPathHeight}[thm] ist die Höhe von \(x\) größer als die Höhe aller
Knoten des bisherigen Weges; \FormulaRefAuto{FinitePathFreshEndpointExtension}[thm] liefert den verlängerten Weg. Ferner ist \(x\neq r\). Der vorletzte
Knoten dieses Weges ist \(a\); nach
\FormulaRefAuto{TreeParentDef}[def] gilt die behauptete Elternrelation.
Umgekehrt ist ein Baumelter der vorletzte Knoten des eindeutigen
Wurzelwegs. Dieser Weg läuft nach
\FormulaRefAuto{B27RankedRootPathForward}[thm] vollständig in Richtung von \(D\). Seine letzte Kante ist also
\((a,x)\in D\).
